\documentclass[a4paper,10pt]{article}
\usepackage[utf8]{inputenc}
\usepackage[T1]{fontenc}
\usepackage[sc]{mathpazo}
\usepackage[scaled=0.88]{helvet}
\usepackage{amsmath,amssymb,mathrsfs,stmaryrd}
\usepackage{microtype}
\usepackage[a4paper,top=1.5cm,bottom=1.5cm,left=1.4cm,right=1.4cm,headheight=12pt,headsep=0.4cm,footskip=0.7cm]{geometry}
\usepackage{xcolor}
\definecolor{Ink}{HTML}{1C2533}\definecolor{Bleu}{HTML}{1F3A5F}\definecolor{BleuMoyen}{HTML}{4A6A8F}
\definecolor{BleuClair}{HTML}{9DB2C9}\definecolor{Fond}{HTML}{F2F4F7}\definecolor{Ocre}{HTML}{B07A26}
\definecolor{Brique}{HTML}{A23B2A}\definecolor{Gris}{HTML}{5E6672}\definecolor{Filet}{HTML}{C5CCD6}
\usepackage{multicol,enumitem,booktabs,array,needspace}
\usepackage[most]{tcolorbox}
\usepackage{tikz,pgfplots}
\pgfplotsset{compat=1.17,trig format plots=rad}
\usetikzlibrary{arrows.meta}
\usepackage{fancyhdr}
\usepackage[colorlinks=true,linkcolor=Bleu,bookmarksopen=true,bookmarksopenlevel=0]{hyperref}

\setlength{\parindent}{0pt}\setlength{\parskip}{0.25em}
\setlength{\columnsep}{0.55cm}\setlength{\columnseprule}{0.3pt}\renewcommand{\columnseprulecolor}{\color{Filet}}
\setlist{itemsep=0.05em,topsep=0.15em,parsep=0pt,leftmargin=1.1em}
\setlist[itemize]{label=\textcolor{BleuMoyen}{\scriptsize\textbullet}}
\setlist[enumerate]{label=\textsf{\scriptsize\color{Bleu}\arabic*.},leftmargin=1.3em}
\linespread{1.02}

\newcommand{\K}{\mathbb{K}}\newcommand{\R}{\mathbb{R}}\newcommand{\C}{\mathbb{C}}\newcommand{\Z}{\mathbb{Z}}
\newcommand{\Mn}[1]{\mathcal{M}_{n}(#1)}\newcommand{\rg}{\operatorname{rg}}\newcommand{\vect}{\operatorname{Vect}}
\newcommand{\Mat}{\operatorname{Mat}}\newcommand{\Id}{\operatorname{Id}}\newcommand{\com}{\operatorname{com}}
\newcommand{\vv}[1]{\overrightarrow{#1}}\newcommand{\ve}[1]{\vec{#1}}
\newcommand{\dt}[1]{\left|\begin{matrix}#1\end{matrix}\right|}
\newcommand{\HP}{{\sffamily\bfseries\scriptsize\textcolor{Brique}{[HP]}}}
\newcommand{\livre}[1]{{\sffamily\scriptsize\color{Gris}[#1]}}

% ---------- titres ----------
\newcommand{\chapitre}[2]{\clearpage\phantomsection\pdfbookmark[0]{#1 --- #2}{chap#1}\markboth{#1 --- #2}{}%
  \noindent\colorbox{Bleu}{\parbox{\dimexpr\linewidth-2\fboxsep}{\color{white}\sffamily\bfseries\Large\strut #1\quad
  {\normalsize\mdseries #2}}}\par\vspace{0.4em}}
\newcommand{\notion}[1]{\par\needspace{9\baselineskip}\vspace{0.5em}\phantomsection\pdfbookmark[1]{#1}{\detokenize{#1}}\noindent{\sffamily\bfseries\color{Bleu}\large #1}\par\nobreak
  \noindent{\color{Ocre}\rule{\linewidth}{0.8pt}}\par\nobreak\vspace{0.1em}}
\newcommand{\sous}[1]{\par\needspace{6\baselineskip}\vspace{0.25em}\noindent{\sffamily\bfseries\small\color{BleuMoyen}#1}\par\nobreak}

% ---------- encadrés ----------
\tcbset{base/.style={enhanced,arc=0pt,boxrule=0pt,left=4pt,right=3pt,top=2pt,bottom=2pt,before skip=0.35em,after skip=0.35em,
  fontupper=\small}}
\newtcolorbox{DEF}[1][]{base,colback=white,borderline west={1.4pt}{0pt}{Bleu},
  before upper={{\sffamily\bfseries\scriptsize\color{Bleu}DÉFINITION\ifx\relax#1\relax\else\ --- #1\fi}\par}}
\newtcolorbox{PROP}[1][]{base,colback=Fond,borderline west={2pt}{0pt}{Bleu},
  before upper={{\sffamily\bfseries\scriptsize\color{Bleu}\ifx\relax#1\relax PROPRIÉTÉ\else #1\fi}\par}}
\newtcolorbox{FORM}[1][]{base,colback=white,colframe=BleuMoyen,boxrule=0.5pt,
  before upper={{\sffamily\bfseries\scriptsize\color{BleuMoyen}\ifx\relax#1\relax FORMULES\else #1\fi}\par}}
\newtcolorbox{METH}[1]{base,colback=white,borderline north={0.8pt}{0pt}{Bleu},borderline south={0.8pt}{0pt}{Bleu},
  before upper={{\sffamily\bfseries\scriptsize\colorbox{Bleu}{\color{white}MÉTHODE}\ \color{Bleu}#1}\par\smallskip}}
\newtcolorbox{PIEGE}{base,colback=Brique!5,borderline west={1.4pt}{0pt}{Brique},
  before upper={{\sffamily\bfseries\scriptsize\color{Brique}PIÈGES}\par}}
\newtcolorbox{ASTU}{base,colback=Ocre!7,borderline west={1.4pt}{0pt}{Ocre},
  before upper={{\sffamily\bfseries\scriptsize\color{Ocre}ASTUCES}\par}}
\newtcolorbox{HPBOX}{base,colback=white,borderline west={1.4pt}{0pt}{Brique},
  before upper={{\sffamily\bfseries\scriptsize\color{Brique}HORS PROGRAMME}\par}}

\pagestyle{fancy}\fancyhf{}
\fancyhead[L]{\scriptsize\sffamily\color{Gris}\leftmark}
\fancyhead[R]{\scriptsize\sffamily\scshape\color{Gris}Le Collé PT --- fiche Algèbre linéaire et réduction}
\fancyfoot[C]{\scriptsize\sffamily\thepage}
\renewcommand{\headrulewidth}{0.3pt}

\pgfplotsset{mini/.style={axis lines=middle,axis line style={Filet},xtick=\empty,ytick=\empty,width=3.2cm,height=3.2cm,
  xmin=-1.15,xmax=1.15,ymin=-1.15,ymax=1.15,title style={font=\scriptsize\sffamily,yshift=-0.15cm},samples=80}}


\newcommand{\Ker}{\operatorname{Ker}}\renewcommand{\Im}{\operatorname{Im}}\newcommand{\tr}{\operatorname{tr}}\newcommand{\Sp}{\operatorname{Sp}}
\newcommand{\GLn}[1]{\mathrm{GL}_n(#1)}
\hypersetup{pdftitle={Le Collé — fiche de référence : algèbre linéaire et réduction}}
\begin{document}
\thispagestyle{empty}
\noindent\colorbox{Bleu}{\parbox{\dimexpr\linewidth-2\fboxsep}{\vspace{0.4em}
{\color{white}\sffamily\scriptsize\textls[200]{LE COLLÉ\ \ ·\ \ MATHÉMATIQUES PT\ \ ·\ \ CHAPITRE 2}}\par\smallskip
{\color{white}\sffamily\bfseries\LARGE Fiche de référence --- Algèbre linéaire et réduction}\par\smallskip
{\color{white!85!Bleu}\sffamily compléments, éléments propres, diagonalisation, trigonalisation, pièges et résultats du TD}\vspace{0.4em}}}

\vspace{0.6em}
{\small Le livre sert à apprendre, la présentation à comprendre, cette fiche à travailler et réviser. Mêmes numéros de parties et
mêmes notations que la présentation et le livre ; \livre{th.~1.27} renvoie au livre, \livre{TD~13} à l'exercice~13 du TD complet.
$\K=\R$ ou $\C$, $f\in\mathcal L(E)$, $\dim E=n$, $A\in\Mn{\K}$.}

\vspace{0.4em}
{\small\sffamily\noindent\begin{tabular}{@{}p{0.6\linewidth}@{}}
1.1 Matrices semblables, sommes directes, trace, hyperplans \dotfill\pageref{n11}\\
1.2 Sous-espaces stables \dotfill\pageref{n12}\\
1.3 Éléments propres \dotfill\pageref{n13}\\
1.4 Polynôme caractéristique \dotfill\pageref{n14}\\
1.5 Diagonalisation \dotfill\pageref{n15}\\
1.6 Applications \dotfill\pageref{n16}\\
1.7 Trigonalisation \dotfill\pageref{n17}\\
1.8 Quelle méthode ? Pièges ; résultats du TD \dotfill\pageref{n18}\\
\end{tabular}}

\vspace{0.6em}
\pdfbookmark[0]{Chapitre 2 --- Algèbre linéaire et réduction}{chap2}\markboth{Chapitre 2 --- Algèbre linéaire et réduction}{}
\begin{multicols}{2}\raggedcolumns
% ======================================================================
\notion{1.1 Semblables, sommes directes, trace, hyperplans}\label{n11}

\begin{FORM}[CHANGEMENT DE BASE \livre{prop.~1.1}]
$P$ = matrice de passage de $\mathcal B$ à $\mathcal B'$ (colonnes : $\mathcal B'$ dans $\mathcal B$) :
\[\Mat_{\mathcal B'}(f)=P^{-1}\Mat_{\mathcal B}(f)\,P.\]
$A$, $B$ \emph{semblables} : $B=P^{-1}AP$. Mêmes rang, $\det$, $\tr$, $\chi$, valeurs propres (avec multiplicités).
\end{FORM}

\begin{PROP}[SOMME DIRECTE \livre{prop.~1.4, th.~1.5}]
$F_1+\dots+F_p$ directe $\iff$ ($x_1+\dots+x_p=0$, $x_i\in F_i\Rightarrow$ tous nuls).\\
$p=2$ : $F\cap G=\{0\}$.\\
Base adaptée : concaténation des bases ; $\dim\bigoplus F_i=\sum\dim F_i$.\\
$E=F\oplus G\iff F\cap G=\{0\}$ et $\dim F+\dim G=n$.
\end{PROP}
\begin{PIEGE}
$p\geq3$ : intersections deux à deux nulles $\not\Rightarrow$ directe ($\vect(1,0)$, $\vect(0,1)$, $\vect(1,1)$).
\end{PIEGE}

\begin{FORM}[TRACE \livre{prop.~1.7}]
$\tr A=\sum a_{i,i}$, linéaire ; $\tr(AB)=\tr(BA)$ ; $\tr(P^{-1}AP)=\tr A$ ; $\tr f$ indépendante de la base.\\
Projecteur : $\tr p=\rg p$. \ $AB-BA=I_n$ impossible \livre{TD~8}.\\
$\tr(ABC)=\tr(CAB)$ mais $\neq\tr(BAC)$ en général.
\end{FORM}

\begin{PROP}[HYPERPLANS \livre{prop.~1.10}]
Hyperplan : supplémentaire d'une droite, dimension $n-1$ ; $=\Ker\varphi$, $\varphi$ forme linéaire non nulle ; équation
$a_1x_1+\dots+a_nx_n=0$.\\
$u\notin H\Rightarrow E=H\oplus\vect(u)$.\\
$p$ hyperplans : intersection de dimension $\geq n-p$ ; sous-espace de dimension $n-p$ = $p$ équations.
\end{PROP}

% ======================================================================
\notion{1.2 Sous-espaces stables}\label{n12}

\begin{DEF}
$F$ stable : $f(F)\subset F$ ; induit $f_F\in\mathcal L(F)$. Vérifier $f(e_j)\in F$ sur une base de $F$.\\
$f\circ g=g\circ f\Rightarrow\Ker g$, $\Im g$ stables par $f$.
\end{DEF}

\begin{PROP}[BLOCS \livre{th.~1.13}]
Base adaptée à $F$ ($\dim F=p$) : $F$ stable $\iff\Mat(f)=\begin{pmatrix}A&B\\0&C\end{pmatrix}$, $A=\Mat(f_F)$.\\
$E=F\oplus G$, les deux stables $\iff\begin{pmatrix}A&0\\0&C\end{pmatrix}$.\\
$\det f=\det A\det C$, \ $\tr f=\tr A+\tr C$.
\end{PROP}

\begin{FORM}[PROJECTEURS, SYMÉTRIES]
$p^2=p$ : $E=\Im p\oplus\Ker p$, $\Mat=\mathrm{diag}(I_r,0)$.\\
$s^2=\Id$ : $E=\Ker(s-\Id)\oplus\Ker(s+\Id)$, $\Mat=\mathrm{diag}(I_r,-I_{n-r})$.\\
Droite $\vect(u)$ stable $\iff f(u)=\lambda u$ \livre{prop.~1.14}.
\end{FORM}

% ======================================================================
\notion{1.3 Éléments propres}\label{n13}

\begin{DEF}[\livre{déf.~1.15}]
$\lambda$ valeur propre : $\exists x\neq0$, $f(x)=\lambda x$ ; $x$ vecteur propre (\emph{non nul}).\\
$E_\lambda(f)=\Ker(f-\lambda\Id_E)$, stable ; $\neq\{0\}\iff\lambda\in\Sp(f)$. \ $E_0=\Ker f$.\\
Matrice : $AX=\lambda X$, $E_\lambda(A)$, $\Sp(A)$.\\
Dérivation sur $\mathcal C^\infty$ : $E_\lambda=\vect(e^{\lambda t})$ ; $P\mapsto XP'$ : $X^k$, valeur $k$.
\end{DEF}

\begin{PROP}[SOMME DIRECTE \livre{th.~1.18}]
Sous-espaces propres pour des valeurs distinctes : somme directe. Vecteurs propres de valeurs distinctes : famille libre.\\
Au plus $n$ valeurs propres \livre{cor.~1.19}.
\end{PROP}

\begin{METH}{Sous-espace propre \livre{méth.~1.1}}
Résoudre $(A-\lambda I_n)X=0$ ;\quad $\dim E_\lambda=n-\rg(A-\lambda I_n)$.\\
$\lambda\in\Sp(A)\iff\det(A-\lambda I_n)=0$. Le rang d'abord.
\end{METH}
\begin{PIEGE}
$\Sp_\R(A)\subset\Sp_\C(A)$ : la rotation $\begin{pmatrix}0&-1\\1&0\end{pmatrix}$ n'a pas de valeur propre réelle ($\pm i$ dans $\C$).
\end{PIEGE}

% ======================================================================
\notion{1.4 Polynôme caractéristique}\label{n14}

\begin{FORM}[DÉFINITION \livre{prop.~1.21}]
$\chi_A(x)=\det(xI_n-A)$, unitaire, degré $n$ :
\[\chi_A=x^n-\tr(A)x^{n-1}+\dots+(-1)^n\det A.\]
$n=2$ : $\chi_A=x^2-\tr(A)\,x+\det A$.\\
$\Sp(A)$ = racines de $\chi_A$ dans $\K$ ; semblables $\Rightarrow$ même $\chi$ ; triangulaire : valeurs propres = diagonale.
\end{FORM}

\begin{PROP}[MULTIPLICITÉ \livre{th.~1.24}]
$m(\lambda)$ = multiplicité comme racine de $\chi_f$.\quad $1\leq\dim E_\lambda\leq m(\lambda)$.\\
Valeur propre simple $\Rightarrow\dim E_\lambda=1$.\\
Preuve : base adaptée à $E_\lambda$ : $\begin{pmatrix}\lambda I_d&B\\0&C\end{pmatrix}$, $\chi_f=(x-\lambda)^d\chi_C$.
\end{PROP}

\begin{METH}{Calculer $\chi_A$ \livre{méth.~1.2}}
Structure d'abord (lignes de somme constante, colonne à faire apparaître un facteur, ligne avec des zéros) ; \emph{factoriser au
fur et à mesure}.\\
Contrôles : $\chi_A(0)=(-1)^n\det A$ ; $\sum\lambda_i=\tr A$ ; $\prod\lambda_i=\det A$.\\
Ex. : $\begin{pmatrix}0&1&1\\1&0&1\\1&1&0\end{pmatrix}$ : $(x-2)(x+1)^2$.
\end{METH}
\begin{PIEGE}
$\det(A-xI_n)=(-1)^n\chi_A(x)$ : mêmes racines, mais ce n'est pas $\chi_A$ si $n$ est impair.
\end{PIEGE}

% ======================================================================
\notion{1.5 Diagonalisation}\label{n15}

\begin{DEF}
$f$ diagonalisable : base de vecteurs propres. $A$ diagonalisable : $A=PDP^{-1}$, colonnes de $P$ = vecteurs propres, $D$ = valeurs
propres \emph{dans le même ordre}.
\end{DEF}

\begin{PROP}[CRITÈRES \livre{th.~1.27}]
$f$ diagonalisable $\iff E=\bigoplus E_\lambda\iff\sum\dim E_\lambda=n$\\
$\iff\chi_f$ scindé et $\dim E_\lambda=m(\lambda)$ pour tout $\lambda$.\\
$n$ valeurs propres distinctes $\Rightarrow$ diagonalisable \livre{cor.~1.28} (réciproque fausse : $I_n$).\\
Une seule valeur propre $\lambda$ et diagonalisable $\Rightarrow A=\lambda I_n$.\\
Projecteurs, symétries : diagonalisables.
\end{PROP}

\begin{METH}{Diagonaliser ($n=2,3$) \livre{méth.~1.3}}
\begin{enumerate}
\item $\chi_A$ factorisé ; non scindé sur $\K$ $\Rightarrow$ non.
\item Racines simples : oui. Multiples : $\dim E_\lambda=n-\rg(A-\lambda I_n)$ vs $m(\lambda)$.
\item Bases des $E_\lambda$ en colonnes : $P$ ; $D$ dans le même ordre.
\item Contrôle : $AP=PD$ colonne par colonne.
\end{enumerate}
\end{METH}

\begin{FORM}[EXEMPLES DU COURS]
$\begin{pmatrix}5&1&-1\\2&4&-2\\1&-1&3\end{pmatrix}$ : $2,4,6$ ; $E_2=\vect(0,1,1)$, $E_4=\vect(1,0,1)$, $E_6=\vect(1,1,0)$.\\
$\begin{pmatrix}3&-1&1\\0&2&0\\1&-1&3\end{pmatrix}$ : $(x-2)^2(x-4)$, $\rg(B-2I)=1$ : diagonalisable.\\
$\begin{pmatrix}2&1\\-1&4\end{pmatrix}$ : $(x-3)^2$, $\dim E_3=1$ : non.
\end{FORM}

% ======================================================================
\notion{1.6 Applications}\label{n16}

\begin{FORM}[PUISSANCES \livre{prop.~1.33}]
$A=PDP^{-1}\Rightarrow A^k=PD^kP^{-1}$, $D^k=\mathrm{diag}(\lambda_i^k)$ ; $k\in\Z$ si $A$ inversible.\\
$\begin{pmatrix}1&2\\2&1\end{pmatrix}^k=\frac12\begin{pmatrix}3^k+(-1)^k&3^k-(-1)^k\\3^k-(-1)^k&3^k+(-1)^k\end{pmatrix}$.
\end{FORM}

\begin{METH}{Suites couplées, récurrences \livre{méth.~1.4}}
$X_{k+1}=AX_k\Rightarrow X_k=A^kX_0$. Coordonnées dans la base propre : $Y=P^{-1}X$, $Y_{k+1}=DY_k$, suites géométriques.\\
Ordre 3 : $X_k=(u_k,u_{k+1},u_{k+2})$, matrice compagnon ; racines distinctes $\Rightarrow u_k=\sum\alpha_i\lambda_i^k$.
\end{METH}

\begin{METH}{Équations matricielles}
$D$ diagonale à coefficients distincts : $MD=DM\iff M$ diagonale \livre{prop.~1.35}.\\
$A=PDP^{-1}$ : poser $N=P^{-1}MP$. $M^2=A$ ($3$ valeurs propres $>0$ distinctes) : $8$ solutions.\\
$AB=BA$ : les $E_\lambda(A)$ sont stables par $B$.
\end{METH}
\begin{HPBOX}
Polynômes d'endomorphismes (polynôme annulateur, minimal) ; systèmes $X'=AX$ (se traiteraient avec $Y=P^{-1}X$) ; techniques générales de trigonalisation.
\end{HPBOX}

% ======================================================================
\notion{1.7 Trigonalisation}\label{n17}

\begin{PROP}[\livre{th.~1.37, prop.~1.38}]
Trigonalisable $\iff\chi$ scindé (admis). Toute $A\in\Mn{\C}$ est trigonalisable.\\
$\chi_A=\prod(x-\lambda_i)$ (avec multiplicités) :
\[\tr A=\sum\lambda_i,\qquad\det A=\prod\lambda_i.\]
\end{PROP}

\begin{ASTU}
$2\times2$ : $\tr$ et $\det$ donnent les deux valeurs propres.\\
$\rg A=1$ : $0$ d'ordre $\geq n-1$, et la dernière est $\tr A$ ; diagonalisable $\iff\tr A\neq0$ \livre{TD~27}.\\
$A^p=0$ : $\Sp_\C(A)=\{0\}$, $\chi_A=x^n$, $\tr A=0$ ; diagonalisable $\iff A=0$.
\end{ASTU}

% ======================================================================
\notion{1.8 Quelle méthode ? Pièges ; résultats du TD}\label{n18}

{\small\renewcommand{\arraystretch}{1.12}
\begin{tabular}{@{}>{\raggedright\arraybackslash}p{0.34\linewidth}>{\raggedright\arraybackslash}p{0.6\linewidth}@{}}
\hline
\textbf{Question} & \textbf{Méthode}\\\hline
$E=F\oplus G$ ? & intersection nulle + dimensions\\
$F$ stable ? & images d'une base de $F$ ; blocs\\
$\lambda$ valeur propre ? & $\rg(A-\lambda I)<n$\\
valeurs propres & $\chi_A$ factorisé ; $\tr$, $\det$\\
diagonalisable ? & distinctes ; sinon $\dim E_\lambda=m(\lambda)$\\
$A^k$, suites & $PD^kP^{-1}$ ; base propre\\
polynômes, fonctions & $f(x)=\lambda x$ directement\\\hline
\end{tabular}}

\begin{PIEGE}
Vecteur propre \emph{non nul}. Oublier de vérifier $\chi$ scindé sur $\R$. Valeur double $\not\Rightarrow$ non diagonalisable.
Vecteurs propres en \emph{colonnes}, dans l'ordre de $D$. $A=PDP^{-1}$, pas $P^{-1}DP$.
\end{PIEGE}

{\small
\renewcommand{\arraystretch}{1.12}
\begin{tabular}{@{}>{\sffamily\scriptsize\bfseries\color{Bleu}}l>{\raggedright\arraybackslash}p{0.84\linewidth}@{}}
TD2 & $E_2=\vect(1,2)$, $E_{-3}=\vect(1,1)$ : $A=\begin{pmatrix}-8&5\\-10&7\end{pmatrix}$.\\
TD7 & $\begin{pmatrix}1&-1&-1\\-1&1&-1\\-1&-1&1\end{pmatrix}$ : $(x+1)(x-2)^2$, diagonalisable.\\
TD11 & $(X+1)P'$ : $(X+1)^k$, valeur $k$.\\
TD12 & $\begin{pmatrix}1&a&1\\0&1&1\\0&0&2\end{pmatrix}$ diagonalisable $\iff a=0$.\\
TD14 & anti-diagonale : symétrie, $\dim E_1=\lceil\frac n2\rceil$.\\
TD15 & $M\mapsto M^T$ : $\mathcal S_n\oplus\mathcal A_n$, $\tr=n$.\\
TD17 & transition : $A^k\to\frac1{p+q}\begin{pmatrix}q&q\\p&p\end{pmatrix}$.\\
TD18 & $J_n$ : $\{0,n\}$ ; $\det(aI+bJ)=a^{n-1}(a+nb)$.\\
TD20 & circulante : $\det=\prod_{\omega^3=1}(a+b\omega+c\omega^2)$.\\
TD21 & $nXP-(X^2-1)P'$ : $(X-1)^k(X+1)^{n-k}$, valeur $n-2k$.\\
TD23 & $\Sp(AB)=\Sp(BA)$.\\
TD24 & $(f-3\Id)(f+2\Id)=0\Rightarrow$ diagonalisable.\\
TD26 & $T\sim\begin{pmatrix}1&0&0\\0&2&1\\0&0&2\end{pmatrix}$, non diagonalisable.\\
TD28 & $u_k=1+(-1)^k+2^k$.
\end{tabular}}

\end{multicols}
\end{document}
